% Generated draft: requires mathematical and source review.
\documentclass[11pt]{article}
\usepackage[margin=25mm]{geometry}
\usepackage{fontspec}
\setmainfont{Latin Modern Roman}
\usepackage{amsmath,amssymb,mathtools,unicode-math}
\setmathfont{Latin Modern Math}
\usepackage{enumitem,needspace,xurl,hyperref}
\setlength{\parindent}{0pt}
\setlength{\parskip}{5pt}
\setlength{\emergencystretch}{3em}
\newcommand{\N}{\mathbb{N}}
\newcommand{\Z}{\mathbb{Z}}
\newcommand{\Q}{\mathbb{Q}}
\newcommand{\R}{\mathbb{R}}
\newcommand{\C}{\mathbb{C}}
\newcommand{\F}{\mathbb{F}}
\newcommand{\A}{\mathbb{A}}
\newcommand{\PP}{\mathbb{P}}
\newcommand{\E}{\mathbb{E}}
\newcommand{\eps}{\varepsilon}
\newcommand{\dd}{\,\mathrm{d}}
\begin{document}
\section*{UnsolvedMath ID: 30004786}
\textbf{Source problem code:} catalogue entry
\subsection{Definitions and mathematical statement}
Write \(L(s,\sigma,\rho)\) for the automorphic \(L\)-function named in the supplied record, with complex variable \(s\). The record does not specify the underlying global field, the group, the class of automorphic representations \(\sigma\), the representation \(\rho\), or the normalization of this \(L\)-function.

Let \(\mathcal P_{\sigma,\rho}\) denote the conjectural \((\sigma,\rho)\)-Poisson summation formula referred to in the source. This notation is a placeholder: the supplied record does not state that formula, its test-function space, its transform, its convergence conditions, or any correction terms. Likewise, let \(\mathcal{FE}_{\sigma,\rho}\) denote the intended functional equation, whose precise form is not supplied.

The requested implication is
\[
\mathcal P_{\sigma,\rho}
\quad\Longrightarrow\quad
\left\{
\begin{array}{l}
L(s,\sigma,\rho)\text{ admits meromorphic continuation},\\
L(s,\sigma,\rho)\text{ satisfies }\mathcal{FE}_{\sigma,\rho}.
\end{array}
\right.
\]
It is intended to apply to the pairs \((\sigma,\rho)\) covered by the conjecture. However, neither that quantified class nor the initial domain and target domain of continuation are specified in the supplied record. The record also does not distinguish a completed \(L\)-function from a partial or finite-part \(L\)-function, or specify the dual data, normalization, and factors in the functional equation. Consequently, this is a conditional statement schema requiring those data before it becomes a fully specified mathematical assertion.

The supplied background additionally discusses an established \(\mathrm{GL}_n\) identity, refined Poisson formulas for selected Braverman--Kazhdan spaces, and an expected direct proof for all split classical groups with the standard representation. These are supplied contextual qualifications; the record does not explain whether the last scope restricts the displayed implication, concerns proving its antecedent, or describes a separate extension.
\subsection{Short English statement}
Show that the conjectural \((\sigma,\rho)\)-Poisson summation formula yields meromorphic continuation and the appropriate functional equation for \(L(s,\sigma,\rho)\). The supplied record leaves the admissible data and both formulas unspecified, so the statement remains a conditional schema requiring human mathematical review.
\subsection{Sources}
\begin{itemize}
\item Local UnsolvedMath registry: supplied record with permanent ID \(30004786\), title \emph{Automorphic L-Functions from Sigma--Rho Poisson Summation}, collection \emph{Oberwolfach Reports Open Problems}, exact source code \texttt{catalogue entry}. The record gives proposed year 2021 and names no proposer.
\item Supplied source citation: \emph{Automorphic Forms, Geometry and Arithmetic} (2022), \emph{Oberwolfach Reports}, pp.~2089--2156, DOI \texttt{10.4171/owr/2021/39}. \url{https://doi.org/10.4171/owr/2021/39}. The literature-source list labels this reference \emph{Oberwolfach Report 39/2021}.
\item Supplied literature reference: \emph{A refined Poisson summation formula}. \url{https://arxiv.org/abs/1707.06091}.
\item Supplied literature reference: \emph{Harmonic analysis on certain spherical varieties}. \url{https://arxiv.org/abs/2103.10261}.
\item Literature assessment dated exactly \texttt{2026-08-22}: the supplied curation labels the problem open, states that the \(\mathrm{GL}_n\) identity is established and refined Poisson formulas are known for selected Braverman--Kazhdan spaces, and describes the proposed direct proof for all split classical groups with the standard representation as an expected extension rather than a general theorem. This is supplied dated curation derived from the OWR record, not an independently verified literature assessment.
\item The registry labels the statement \texttt{corrected\_verified} and reports that its standalone formulation and hypotheses were checked against the official Oberwolfach report. This verification claim is supplied metadata; the report and cited papers have not been independently consulted for this draft.
\end{itemize}
\end{document}
